\section{Discussion}\label{sec:discussion}

\subsection{Answer to the question}
\begin{itemize}[itemsep=1pt]
\item \textbf{Can steering do it?} Yes. At the frequency (held exactly in the design
  model, to within $\nHeldTol\%$ on refined bases and an independent finite-element model) and mass of the unidirectional plate, a manufacturable steered tow reduces the frequency-normalized
  hardening by $\nFTenRed\%$. No uniform layup can reach that frequency at a lower
  hardening, and steering undercuts the best stacking sequences at every higher
  frequency we test.
\item \textbf{Why?} The frequency and the hardening read the same rotated stiffness
  through two strain fields whose principal directions turn differently over the plate,
  and an optimal design must satisfy the trade-off between them pointwise
  (Proposition~\ref{prop:pointwise}).
\item \textbf{What bounds it?} The material anisotropy ($U_2,U_3$), and the tow radius
  relative to the length scale of the fields.
\item \textbf{How does it compare?} Steering is the stronger channel here. It is also
  the one that does not require changing the bill of materials.
\end{itemize}

\begin{table}[h]
\centering\footnotesize
\caption{Status of the predictions.}\label{tab:verdict}
\begin{tabular}{@{}p{0.39\linewidth}p{0.39\linewidth}p{0.13\linewidth}@{}}
\toprule
prediction & evidence & verdict\\ \midrule
P1 isotropic material has no lever & Cor.~\ref{cor:iso}; gradient $\equiv0$ to $10^{-10}$; Fig.~\ref{fig:limits}a & holds (exact)\\
P2 uniform angle cannot change a single-ply beam; steering can, at the inflection points & SSM ratio $\nBUnif$; Tab.~\ref{tab:beam}; Fig.~\ref{fig:beam} & holds\\
P3 steered plate reaches $(\omega_1,J)$ no uniform layup reaches & Fig.~\ref{fig:levers}; $\nSnum$ stacks & holds\\
P4 optimized field follows preferred directions; $\theta^\star$ captures most of the gain & Fig.~\ref{fig:mechanism}; $\theta^\star$: $\nStarJ$ at $\nStarF\,\omega_{1,\UD}$; optimum departs by $\nStarDevHigh^\circ$ RMS where energies are large & holds partly\\
P5 lever grows with anisotropy, saturates at the field length scale & Fig.~\ref{fig:limits} & holds\\
consequence: $a_c$ widens as $J^{-1/2}$ & $\nCacRatio$ vs.\ $\nCacTheory$; time-domain sweeps & holds\\
\bottomrule
\end{tabular}
\end{table}

\subsection{What is new}
\begin{itemize}[itemsep=1pt]
\item \textbf{Design variable.} Reduced-order optimization of nonlinear dynamics
  \cite{pozzi2024lsm,pozzi2025,pozzi2025smo,liang2025,xiao2026} has used shape, size, density or other system parameters. We use a
  manufacturable fiber-orientation field.
\item \textbf{Objective.} Variable-stiffness composites \cite{gurdal,ijsselmuiden} have been
  steered for linear objectives and for postbuckling \cite{wu2013,raju2018}, and for the
  linear frequency with the nonlinear response examined afterwards
  \cite{farsadi2021,farsadi2025}. We steer them for the hardening itself, at fixed linear
  frequency.
\item \textbf{Mechanism, not only an optimum.} The paper's main contribution is the
  mechanism: a local Fourier theory that predicts when steering must help, where the
  fibers should go, and what limits the gain, tested prediction by prediction.
\item \textbf{Analysis versus design.} Nonlinear vibrations of variable-stiffness plates
  have been computed for prescribed fiber paths
  \cite{ribeiro2012,ribeiroakhavan2012,akhavan2015,akhavan2017,houmat2013,yan2021}. Here
  the path is designed at a held frequency, and the reason it works is identified.
\item \textbf{One-dimensional result.} On a clamped beam the operative lever is not the
  bending--extension coupling but the harmonic-versus-curvature-weighted mean of
  Proposition~\ref{prop:beam}, which requires no coupling at all.
\end{itemize}

\subsection{Limitations and scope}
\begin{itemize}[itemsep=2pt]
\item \textbf{Geometry and support.} Results are for a simply supported, immovable-edge
  rectangular plate. The mechanism is general, but the numbers depend on where the
  curvature and the stretch live, which other supports change.
\item \textbf{Constant thickness.} We assume constant thickness and mass. Real AFP
  steering produces tow gaps or overlaps whose thickness and stiffness changes are not
  modeled. Including them, through a thickness field tied to $\theta(\bm x)$, is the main
  step toward a buildable design.
\item \textbf{Amplitude range.} The design quantity is the small-amplitude cubic
  coefficient. The cubic SSM reproduces the full-order EPMC detuning to within $\nEPMCAgree\%$
  over the amplitudes EPMC reached ($a/h\le\nEPMCmaxUD$ for UD, $\le\nEPMCmaxOpt$ for the steered design),
  and the jump-onset amplitudes reported lie inside these ranges. Higher
  amplitudes need the quintic SSM or EPMC (Appendix~\ref{app:calib}).
\item \textbf{Single mode, no internal resonance.} On the full multi-mode plate the steered design's lowest
  frequencies are in ratios $\nModeRatios$ (\texttt{mode\_\allowbreak ratios.py}). The nearest
  commensurability ($3{:}1$ with the third mode) is detuned by $\nDetThreeOne\%$, and that mode
  is antisymmetric about the center. The $r=6$ mode time-domain model contains these modes
  and would show an internal-resonance exchange; its sweeps follow the single-mode forced
  response instead (Fig.~\ref{fig:confirm}b,c).
\item \textbf{Optimality.} The optimizer is local (multi-start SLSQP). The pointwise
  ceiling of Sec.~\ref{sec:limits} bounds what any single-ply field can do. Stacks have eight
  plies (four independent angles), optimized continuously at each held frequency; more
  plies or a lamination-parameter search could lower the uniform optimum further; the
  steered values are at most $\nSteerVsUniformMax\%$ of the present one. The UD floor at
  $\omega_{1,\UD}$ holds for every uniform laminate (Corollary~\ref{cor:udfloor}).
\item \textbf{Damping.} Damping is fixed ($2\%$ of critical, mass-proportional) and not designed.
\item \textbf{Idealized laminate.} The plate designs are a single $2$\,mm steered ply,
  i.e.\ an unbalanced laminate, and tow gaps and overlaps are ignored. The beam design
  with no tow limit turns to about $\nBeamMaxAngle^\circ$ in narrow zones of the $20$\,mm-wide strip.
\end{itemize}

\subsection{Outlook}
The same reduction and adjoint carry over directly to three extensions:
\begin{itemize}[itemsep=1pt]
\item \emph{prescribing} the backbone (a target frequency--amplitude law) instead of
  minimizing the hardening;
\item adding a density field, so that topology and tow path are designed together;
\item steering stability boundaries, such as supersonic panel flutter, whose
  nonsymmetric aeroelastic operator enters the same first-order pencil; the flutter and
  post-flutter response of plates with prescribed curvilinear fibers has been analyzed
  \cite{akhavan2018,akhavan2020,farsadi2019}, but not designed.
\end{itemize}
The mechanism of Sec.~\ref{sec:theory} predicts where each extension should help: wherever
two functionals of the same stiffness are weighted by different strain fields.

\section{Conclusions}\label{sec:conclusions}
\begin{enumerate}[itemsep=2pt]
\item The frequency and the geometric hardening of a laminated panel integrate the same
  rotated stiffness against the curvature and the relaxed stretch of the fundamental.
  Each density depends on the local fiber angle only through $U_2$, $U_3$ and the local
  principal directions (Proposition~\ref{prop:fourier}).
\item A uniform single angle cannot change the frequency-normalized hardening of a
  single-ply beam. Symmetric steering cuts it by up to $\nBInfHeldRed\%$ with the
  frequency within $\nBeamFreqBand\%$ of the UD value, by placing compliant fibers at the inflection points. In our comparison
  bending--extension coupling was a minor lever ($\nBBred\%$ on its own drift objective,
  against $\nBNineNinetyRed\%$ for symmetric steering at the same $900^\circ$/m tow limit and
  frequency floor); a richer coupled parametrization might do better.
\item On a simply supported plate, a curvature-feasible steered field reduces the
  frequency-normalized hardening by $\nFTenRed\%$ at the unidirectional frequency
  (held exactly in the design model, to within $\nHeldTol\%$ on refined bases and an independent finite-element model) and mass. This point is unreachable by any uniform layup, and the steered
  designs undercut the best stacking sequences at every frequency tested.
\item A pointwise design from the UD fields alone captures most of the gain; the
  optimized field shares its large-scale pattern but not its detail. Steering lowers the restrained stretch energy $k_{nl}$ to $\nDknl$ of its UD
  value; the relaxation falls less ($\nDrelax$), so its share of the stretch energy rises
  ($\nRelaxUD\%\to\nRelaxOpt\%$) and $\gamma$ falls further than $k_{nl}$.
\item The lever vanishes for isotropic material, grows monotonically with anisotropy, and
  saturates once the minimum tow radius is below about $\SI{20}{mm}$.
\item The full-order model confirms the reduction (SSM, EPMC, time integration), and the
  jump-free amplitude widens by $\nCacRatio\times$, in agreement with
  $a_c\propto J^{-1/2}$.
\end{enumerate}

\paragraph{Data and code availability.} The code, result files and the scripts that generate
every number, table and figure (Appendix~\ref{app:repro}) are in the repository
\url{https://github.com/jtempelman95/NLTopt} (branch \texttt{paper-revision}; see \texttt{docs/paper/README\_REPRODUCE.md}). The SSM, EPMC and time-integration scripts build on the author's
harmonic-balance package PyHB, which is not public and is available on request; the design
evaluator and the regeneration of every number from the committed result files do not need it.

\paragraph{Competing interests.} The author declares no competing interests.
